\documentclass[french,letterpaper, 12pt]{article}




\usepackage[top = 1.6cm, left = 2cm, right = 2cm ]{geometry}
\usepackage[pdftex]{graphicx}
\usepackage{soulutf8}
\usepackage{amsmath}
\usepackage{tikz}
\usepackage[utf8]{inputenc}
\usepackage{longtable}
\usepackage[T1]{fontenc}
\usepackage{epigraph}
\usepackage{fancyhdr}
\usepackage{float}
\usepackage{xcolor}
\usepackage{eurosym}
\usepackage{calc}
\usepackage{hyperref}
\usepackage{multirow}
\usepackage{caption}
\usepackage[Algorithme]{algorithm}
\usepackage{algorithmic}
\usepackage{enumerate}
\usepackage[french]{babel}
\usepackage{tcolorbox}
\usepackage{multicol}
\usepackage{etoolbox,refcount}
\usepackage{listings}
\usepackage{amssymb}
\usepackage{subcaption}
\usepackage[standard,thref, framed, hyperref,standard,thmmarks]{ntheorem}
\usepackage{lmodern}
\usepackage{thmtools}


%
%%%%%%% Sub librairies

\usetikzlibrary{arrows,automata}
\usetikzlibrary{decorations.pathreplacing,shapes,arrows,positioning}
\usetikzlibrary{calc} 
\usetikzlibrary{positioning}
\usetikzlibrary{snakes}
\usetikzlibrary{shapes,fit}


%
%%%%% Custom commands
%



\lstset{
	keywordstyle=\color{red},
	basicstyle=\scriptsize\ttfamily,
	commentstyle=\ttfamily\itshape\color{gray},
	stringstyle=\ttfamily,
	showstringspaces=false,
	breaklines=true,
	frameround=ffff,
	rulecolor=\color{black}
}

%
\def\changemargin#1#2{\list{}{\rightmargin#2\leftmargin#1}\item[]}
\let\endchangemargin=\endlist 
%
\newcommand{\newlinealinea}{
	~\\ \hspace*{0.5cm}}
%
\newcommand{\alinea}{
	\hspace*{0.5cm}}
%
\newcommand{\alinealong}{
	\hspace*{1.1cm}}
%
\newcommand{\alignparagraph}{
	\hspace*{0.6cm}}
%
\newcommand{\red}[1]{
	\textcolor{red}{#1}}
%
\newcommand{\green}[1]{
	\textcolor{green}{#1}}
%
\newcommand{\point}{$\bullet\ $}
%
\makeatletter
\newcommand*{\whiten}[1]{\llap{\textcolor{white}{{\the\SOUL@token}}\hspace{#1pt}}}
\newcommand{\myul}[1]{
	\underline{\smash{#1}}
}
\makeatother
%
\setlength{\fboxsep}{2pt}
%
\DeclareMathOperator*{\argmax}{\arg\!\max}
%
%
%%%%% Custom text
%
%
\makeatletter
\@addtoreset{section}{part}
\makeatother  
%
\renewcommand*\sfdefault{phv}
\renewcommand*\rmdefault{ppl}
%
\renewcommand\epigraphflush{flushright}
\renewcommand\epigraphsize{\normalsize}
\setlength\epigraphwidth{0.7\textwidth}
%
\definecolor{titlepagecolor}{cmyk}{0.24,0.92,0.78,0.25}
\definecolor{red}{cmyk}{0, 0.91, 0.91, 0.20}
%
\DeclareFixedFont{\titlefont}{T1}{phv}{\seriesdefault}{n}{0.375in}
%
%
%%%%% Header
%
%
\pagestyle{fancy}
\lhead{\student}
\rhead{\grade}
\cfoot{\thepage}
%
%
%%%%% Title page. The following code is borrowed from: 
%%%%%       http://tex.stackexchange.com/a/86310/10898
%
%
\newcommand\titlepagedecoration{%
	\begin{tikzpicture}[remember picture,overlay,shorten >= -10pt]
	
	\coordinate (aux1) at ([yshift=-70pt]current page.north east);
	\coordinate (aux2) at ([yshift=-460pt]current page.north east);
	\coordinate (aux3) at ([xshift=-6cm]current page.north east);
	\coordinate (aux4) at ([yshift=-150pt]current page.north east);
	
	\begin{scope}[titlepagecolor!40,line width=12pt,rounded corners=12pt]
	\draw
	(aux1) -- coordinate (a)
	++(225:5) --
	++(-45:5.1) coordinate (b);
	\draw[shorten <= -10pt]
	(aux3) --
	(a) --
	(aux1);
	\draw[opacity=0.6,titlepagecolor,shorten <= -10pt]
	(b) --
	++(225:2.2) --
	++(-45:2.2);
	\end{scope}
	\draw[titlepagecolor,line width=8pt,rounded corners=8pt,shorten <= -10pt]
	(aux4) --
	++(225:0.8) --
	++(-45:0.8);
	\begin{scope}[titlepagecolor!70,line width=6pt,rounded corners=8pt]
	\draw[shorten <= -10pt]
	(aux2) --
	++(225:3) coordinate[pos=0.45] (c) --
	++(-45:3.1);
	\draw
	(aux2) --
	(c) --
	++(135:2.5) --
	++(45:2.5) --
	++(-45:2.5) coordinate[pos=0.3] (d);   
	\draw 
	(d) -- +(45:1);
	\end{scope}
	\end{tikzpicture}%
}
%%%%%%%%%%%%%%%%%%   ALGORITHM %%%%%%%%%%%%%%%%%%%%%

\renewcommand{\algorithmicrequire}{\textbf{Requis:}}
\renewcommand{\algorithmicensure}{\textbf{Promet:}}
\renewcommand{\algorithmicend}{\textbf{fin}}
\renewcommand{\algorithmicif}{\textbf{si}}
\renewcommand{\algorithmicthen}{\textbf{alors}}
\renewcommand{\algorithmicelse}{\textbf{sinon}}
\renewcommand{\algorithmicelsif}{\algorithmicelse\ \algorithmicif}
\renewcommand{\algorithmicendif}{\algorithmicend\ \algorithmicif}
\renewcommand{\algorithmicfor}{\textbf{pour}}
\renewcommand{\algorithmicforall}{\textbf{pour chaque}}
\renewcommand{\algorithmicdo}{\textbf{faire}}
\renewcommand{\algorithmicendfor}{\algorithmicend\ \algorithmicfor}
\renewcommand{\algorithmicwhile}{\textbf{tant que}}
\renewcommand{\algorithmicendwhile}{\algorithmicend\ \algorithmicwhile}
\renewcommand{\algorithmicloop}{\textbf{boucle}}
\renewcommand{\algorithmicendloop}{\algorithmicend\ \algorithmicloop}
\renewcommand{\algorithmicrepeat}{\textbf{répéter}}
\renewcommand{\algorithmicuntil}{\textbf{jusqu'à}}
\renewcommand{\algorithmicprint}{\textbf{afficher}}
\renewcommand{\algorithmicreturn}{\textbf{retourner}}
\renewcommand{\algorithmictrue}{\textbf{vrai}}
\renewcommand{\algorithmicfalse}{\textbf{faux}}

\renewcommand{\algorithmicand}{\textbf{et}}
\renewcommand{\algorithmicor}{\textbf{ou}}
%%%%%%%%%%%%%%%%%%%%%%%%%%%%%%%%%%%%%%%%%%%%%%%%%%%

\newcommand{\todo}[1]{\textcolor{red}{\emph{\textbf{TODO} : #1}}}
%%%%%%%%%%%%%%%%%%%%%%%%%%%%%%%%%%%%%%%%%%%%%%%%%%%%


\newcounter{countitems}
\newcounter{nextitemizecount}
\newcommand{\setupcountitems}{%
	\stepcounter{nextitemizecount}%
	\setcounter{countitems}{0}%
	\preto\item{\stepcounter{countitems}}%
}
\makeatletter
\newcommand{\computecountitems}{%
	\edef\@currentlabel{\number\c@countitems}%
	\label{countitems@\number\numexpr\value{nextitemizecount}-1\relax}%
}
\newcommand{\nextitemizecount}{%
	\getrefnumber{countitems@\number\c@nextitemizecount}%
}
\newcommand{\previtemizecount}{%
	\getrefnumber{countitems@\number\numexpr\value{nextitemizecount}-1\relax}%
}
\makeatother    
\newenvironment{AutoMultiColItemize}{%
	\ifnumcomp{\nextitemizecount}{>}{3}{\begin{multicols}{3}}{}%
		\setupcountitems\begin{itemize}}%
		{\end{itemize}%
		\unskip\computecountitems\ifnumcomp{\previtemizecount}{>}{3}{\end{multicols}}{}}


%%%%%%%%%%%%%%%%%%%%%%%%%%%%%%%%%%%%%%%%%%%%%%%%%%%%

\theoremstyle{plain}


\renewtheorem{theorem}{Théorème}[section]
\renewtheorem{lemma}[theorem]{Lemme}
\renewtheorem{proposition}[theorem]{Proposition}


\theoremstyle{break}

\newtheorem{algo}{Algorithme}[section]

\theoremheaderfont{\normalfont\smallskip\normalsize\itshape\bfseries}
\theorembodyfont{\normalfont\normalsize}



\renewtheorem{proof}{Preuve}[theorem]
\renewtheorem{example}{Exemple}
\renewtheorem{corollary}{Corollaire}[theorem]


\newtheorem{algoproof}{Preuve}[algo]
\newtheorem{complexity}[algoproof]{Complexité}


\newcommand{\hdelta}{\hat{\delta}}
\newcommand{\automaton}{$A=(Q,\Sigma, q_0, \delta, F)$\ }
\renewcommand{\rm}{$R_M$\ }

%%%%%%%%%%%%%%%%%%%%%%%%%%%%%%%%%%%%%%%%%%%%%%%%%%%%

\newcommand{\student}{Benjamin André}
\newcommand{\grade}{MAB2 Sciences Informatiques}
\newcommand{\director}{Véronique Bruyère}
\renewcommand{\title}{Automates}
\renewcommand{\date}{\today}

\begin{document}
	
	\begin{titlepage}
		%
		\noindent
		%
		\newgeometry{bottom = 2cm, top = 2.5cm}
		\begin{center}
			\includegraphics[scale=1.2]{res/UMONS}\\
			\vspace*{0.3cm}
			\includegraphics[scale=0.23]{res/FS_Logo}\\
			\vspace*{2.5cm}
			%
			\titlefont \title \par
			%
		\end{center}
		\vspace*{3cm}
		\hfill
		%
		\begin{minipage}{0.18\linewidth}
			\begin{flushright}
				\rule{0.5pt}{50pt}
			\end{flushright}
		\end{minipage}
		%
		\begin{minipage}{0.8\linewidth}
			\begin{flushleft}
				\textsf{\textbf{Étudiant:}} \student\\
				\textsf{\textbf{Directrice:}}\director\\
				\date
			\end{flushleft}
		\end{minipage}
		%
		\vspace*{\fill}                                                             
		%
		\begin{center}
			Faculté des Sciences $\bullet$ Université de Mons $\bullet$ 
			Place du Parc 20 $\bullet$ B-7000 Mons
		\end{center}
		%
		\titlepagedecoration
		%
	\end{titlepage}
	%
	%
	%%%% Tables des matières
	%
	%
	\newgeometry{top = 3cm, left = 2cm, right = 2cm, bottom=2.5cm}
	\cleardoublepage
	
	\tableofcontents
	\newpage
	
	\section{Introduction}
	Ce document a pour but d'amener à la compréhension de l'algorithme d'Angluin. Pour ce faire, des bases théoriques seront posées dans la section \ref{sec:theorie}. Ensuite, quelques algorithmes utilisés par la méthode d'Angluin sont présentés et analysés dans la section \ref{sec:algorithmes}. Dans la section \ref{sec:preuves}, certaines notions théoriques supplémentaires sont apportées, se reposant sur les différents algorithmes du point précédant. Finalement, dans la section \ref{?}, l'algorithme d'Angluin est expliqué et illustré avec un exemple.
	
	
	\section{Bases théoriques}\label{sec:theorie}	\subsection{Alphabet}
	
	Un alphabet, nommé par convention $\Sigma$ est un ensemble fini et non vide de symboles.
	Voici certains exemples d'alphabets :
	\begin{itemize}
		\item $\Sigma = \{0,1,2,3,4,5,6,7,8,9\}$, l'alphabet des chiffres
		\item $\Sigma = \{a,b,c,...,z,A,B,C,...,Z\}$, l'alphabet latin
		\item $\Sigma = \{0,1\}$ l'alphabet binaire
	\end{itemize}

	\subsection{Mots}
	
	Comme $\Sigma$ est un ensemble, on peut définir $\Sigma^k$, qui donne des k-uples de symboles, appartenant tous à $\Sigma$.
	
	Un mot $w$ de taille $|w|=k$ est un ensemble de symboles provenant de $\Sigma^k$. Dans le cas particulier où $k=0$, on note le mot vide (sans symbole) $w=\epsilon$.
	
	De façon générale, $w$ est un mot sur $\Sigma$ si il existe $k$ tel que $w \in \Sigma^k$. Par convention, les mots sont nommés par une lettre minuscule, souvent $w,x,y,z$. 
	
	L'ensemble de tous ces mots possible sur $\Sigma$ est noté $\Sigma^*$. Cet ensemble est infini.
	

	
	\subsection{Langage}
	
	Un langage $L$ est défini sur un alphabet $\Sigma$. $L$ est un ensemble de mots sur cet alphabet : $L \subseteq \Sigma^*$. Comme $\Sigma^*$ contient une infinité de mots, $L$ est susceptible de ne pas être fini non plus.
	
	Par exemple, par rapport à tous les mots disponibles avec les lettres de l'alphabet latin ($\Sigma^*$), seulement certains font partie de la langue française ($L$).
	
	$L$ peut être défini  :
	\begin{itemize}
		\item en énumérant les mots en faisant partie : $L=\{12,35,42,7,0\}$
		\item via un notation ensembliste : $L=\{0^k1^j|k+j=7\}$ ou $L=\{w|w \text{ est un mot français}\}$
	\end{itemize}
	Dans ces deux cas, $\Sigma$ est souvent implicite.
	
	\subsubsection{Problème}
	Une notion liée aux langages est celle de problème. Ce que nous appelons couramment problème peut être exprimé en terme d'appartenance d'un mot $w$ à un langage $L$ sur un alphabet $\Sigma$.
	
	Par exemple, prenons l'alphabet $\Sigma=\{0,1,2,3,4,5,6,7,8,9\}$. Considérons ensuite le langage $L = \{w | \text{le nombre représenté par } w \text{ est pair}\}$.
	Demander si un mot $w$ appartient à $L$ revient à résoudre le problème sur le test de parité de $w$. (Le cas échéant, soit $w$ ne représente pas un nombre, soit il représente un nombre impair).
		
	\subsection{Expression régulière}
	
	Une expression régulière est une façon efficace de définir un langage. La construction de l'expression se fait de façon inductive.
	
	Le cas de base est l'expression $e = a, a \in \Sigma$. Dès lors $L_e = {a}$ où $L_e$ est le langage donné par l'expression $e$.
	
	Les autres règles sont inductives. Par ordre de priorité :
	\begin{itemize}
		\item $e = (e_0)$ La mise en évidence. Ici, $L_e = \{w|w \in L_{e_0}\}$
		\item $e = e_0a, a \in \Sigma$. La concaténation. Ici, $L_e = \{wa|w \in L_{e_0}\}$
		\item $e = e_0+e_1$. L'union. Ici, $L_e = L_{e_0} \cup L_{e_1}$
		\item $e = e_0^*$. La fermeture. Intuitivement, il s'agit de tous les mots qui peuvent être formé par une concaténation de mots définis par $e_0$, éventuellement aucun. Ici, $L_e = \{w_0w_1...w_k | k \in \mathbb{N}, w_0,w_1,...,w_k \in L_{e_0}\}$
		\item $e = e_0^+$. La fermeture non nulle. Il s'agit de la fermeture mais avec toujours au moins un mot venant du langage défini par $e_0$. Ici, $L_e = \{w_0w_1...w_k | k \in \mathbb{N}^0, w_0,w_1,...,w_k \in L_{e_0}\}$
	\end{itemize}
	
	Par exemple, on peut écrire l'expression $e_B = (1+01)1^*0(0+1)^*$.\\
	
	\begin{minipage}{0.5\linewidth}
		Les mots suivant en feraient partie :
		\begin{itemize}
			\item 10
			\item 010
			\item 0110
			\item 0111110
			\item 0101101
		\end{itemize}
	\end{minipage}
	\begin{minipage}{0.5\linewidth}
		Ceux-ci n'en feraient pas partie
		\begin{itemize}
			\item 00
			\item 1
			\item 01
			\item 0101
			\item 11
		\end{itemize}
	\end{minipage}

	Un langage pouvant être écrit sous la forme d'une expression régulière est appelé langage régulier. Un des conséquences de cette propriété est qu'il peut être représenté par un automate déterministe fini. \todo{a prouver}
	
	
	\subsection{Automate déterministe fini}\label{sub:dfa}
	Soit un ensemble de symboles $\Sigma$. Soient $\Sigma^* = \{ a_1a_2a_3...a_n | a_1,a_2,a_3,...,a_n \in \Sigma \}$, l'ensemble des mots de taille arbitraire qu'il est possible de former à partir de $\Sigma$ et $|w|, w \in \Sigma$ la longueur de $w$, le nombre de symboles utilisés. Si $|w|=0$, on note $w=\epsilon$.
	
	
	Un automate est défini par $A = (Q, \Sigma, q_0, \delta, F)$ où
	\begin{itemize}
		\item $Q$ est un ensemble d'états, différenciés par leur indice $q_1, q_2, ..., q_n$ ou $n = |Q|$.
		\item $\Sigma$ est un ensemble de symboles
		\item $q_0 \in Q$ est l'état initial
		\item $\delta : Q \times \Sigma \rightarrow Q$ est la fonction de transition. A partir d'un état de $Q$, en fonction d'un symbole, elle retourne un nouvel état faisant partie de $Q$.
		\item $F \subseteq Q$ est un ensemble d'état finaux.
	\end{itemize}
	 
	 
	 Exemple :
	 Soient 
	 \begin{itemize}
	 	\item $\Sigma=\{0,1\}$
	 	\item $Q=\{q_a,q_b,q_c,q_d,q_e,q_f,q_g\}$
	 	\item $q_0=q_a$
	 	\item $F=\{q_d\}$
	 \end{itemize}
 	
 	Pour obtenir un automate, il manque la description de $\delta$. Une façon efficace de noter cette fonction de transition est à l'aide d'une table dont les lignes reprennent des éléments de $Q$, les colonnes des symboles de $\Sigma$, et les cases des éléments de $Q$. Ainsi, on peut lire la relation $\delta : Q \times \Sigma \rightarrow Q$ en prenant une ligne et une colonne.
 	De plus, via cette notation, $Q$ et $\Sigma$ sont explicites. En dénotant l'état initial par $\rightarrow$ et les états acceptants par $*$, on obtient une définition complète d'un automate.
 	
 	\begin{figure}[H]
 	\centering
 	\begin{tabular}{|r||c|c|}
 		\hline
 		&0&1\\
 		\hline\hline
 		$\rightarrow q_a$&$q_c$&$q_b$\\\hline
 		$q_b$&$q_d$&$q_f$\\\hline
 		$q_c$&$q_e$&$q_f$\\\hline
 		$q_d$&$q_d$&$q_d$\\\hline
 		$q_e$&$q_e$&$q_e$\\\hline
 		$q_f$&$q_d$&$q_b$\\\hline
 		$q_g$&$q_e$&$q_f$\\\hline
 	\end{tabular}
	\caption{La table de transitions $\delta_B$}
 	\end{figure}
	 L'automate peut aussi être représenté graphiquement : 
	 
	 \begin{figure}[H]
	 	\centering
	 	\begin{tikzpicture}[->,>=stealth',shorten >=1pt,auto,node distance=3cm, semithick, bend angle=10]
	 	
	 	\tikzstyle{every state}=[circle]
	 	
	 	\node[initial,state] (A)                    {$q_a$};
	 	\node[state]         (B) [below right of=A] {$q_b$};
	 	\node[state]         (C) [below left of=A] {$q_c$};
	 	\node[accepting, state]         (D) [below right of=B] {$q_d$};
	 	\node[state]         (E) [below left of=C]       {$q_e$};
	 	\node[state]         (G) [below right of=E]       {$q_g$};
	 	\node[state]         (F) [above right of=G]       {$q_f$};
	 	
	 	\path 	(A) 	edge              node {0} (C)
	 	edge              node {1} (B)
	 	(B) 	edge              node {0} (D)
	 	edge [bend left]  node {1} (F)
	 	(C) 	edge              node {0} (E)
	 	edge              node {1} (F)
	 	(D) 	edge [loop above] node {0,1} (D)
	 	(E) 	edge [loop above] node {0,1} (E)
	 	(F) 	edge              node {0} (D)
	 	edge [bend left]  node {1} (B)
	 	(G) 	edge              node {0} (E)
	 	edge              node {1} (F);
	 	\end{tikzpicture}
	 	\caption{Automate $A_B$, exemple personnel}\label{fig:ab}
	 \end{figure}
	 
	 Le gain de clarté sur la structure se fait par une perte d'efficacité pour suivre les transitions, surtout quand le graphe est complexe et les arcs s'intersectent. De plus, les ensembles $Q$ et $\Sigma$ sont implicites et doivent être définis à part. $q_0$ et $F$ sont respectivement représenté par la flèche entrante et le double cercle.
	 
	 \subsubsection{Chemin}
	 
	 Pour faire le lien avec la notion de langage, il doit exister une façon pour l'automate de représenter quels mots sont acceptés ou non.
	 
	 Pour ce faire, la fonction $\delta$ peut être étendue à un chemin $w$:
	 
	 \begin{itemize}
	 	\item Si $|w| \leq 1$, $\hat{\delta}(x, w) = \delta(x, w)$
	 	\item Sinon, c'est que $w$ peut s'écrire $au, u \in \Sigma^*$. Alors, $\hat{\delta}(x,w) = \hat{\delta}(x,au) = \hat{\delta}(\delta(x,a),u)$
	 \end{itemize}
	 
	 
	 Le langage représenté par un automate $A$ peut alors se définir comme les mots menant de l'état initial a un état acceptant :
	 $$
	 \{w \in \Sigma^* | \hat{\delta}(q_O,w) \in F_A\}
	 $$
	 
	 
	 
	 
	\section{Algorithmes}\label{sec:algorithmes}Cette section repose sur les bases théoriques, principalement sur l'automate et la relation $R_M$, ainsi que le théorème de Myhill-Nérode, pour donner l'algorithme de remplissage de table (Table Filling Algorithme). A partir de celui-ci, les questions de minimisation, équivalence et constructions d'automates sont explorées.


\subsection{Table Filling Algorithm}\label{ss:tfa}

Le \emph{Table Filling Algorithm} permet, pour un automate, de déterminer quels états sont équivalents. Il repose sur la définition de la relation $R_M$ définie en \ref{ss:rm}

\subsubsection{Construction de la table}

L'idée est de construire, par induction, une table nous disant pour chaque paire d'état si ceux-ci sont équivalents où non, suivant la relation $R_M$ définie précédemment:

\textbf{Cas de base :} Si $p$ est un état final et que $q$ ne l'est pas, alors la paire $\{p,q\}$ est différentiable.

\textbf{Induction : } Soient $p,q$ des états tels qu'il existe un symbole $a$ qui donne $\delta(p,a)=r$ et $\delta(q,a)=s$. Si $r$ et $s$ sont différentiables, alors $p$ et $q$ le sont aussi. En effet, il existe un mot témoin $w$ qui permet de différencier $r$ et $s$. Alors le mot $aw$ permet de différencier $p$ et $q$.

\begin{theorem}
	Si deux états ne sont pas distingués par l'algorithme de remplissage de table, les états sont équivalents.
\end{theorem}

\begin{proof}
	Considérons un automate déterministe fini quelconque $A = (Q, \Sigma, \delta, q_0, F)$, et faisons une preuve par l'absurde.
	
	Supposons qu'il existe une paire d'états $\{p,q\}$ tels que :
	\begin{enumerate}
		\item $p$ et $q$ ne sont pas distingués par l'algorithme de remplissage de table
		\item Les états ne sont pas équivalents, c'est à dire différentiables.
	\end{enumerate}
	
	L'hypothèse deux implique qu'il existe un mot $w \in \Sigma^*$ tel que de $\hat{\delta}(p,w)$ et $\hat{\delta}(q,w)$ un et un seul soit un état final.
	
	Une telle paire est une \emph{mauvaise paire}. Si il y a des mauvaises paires, chacune distinguée par un mot témoin, il doit exister un paire distinguée par le mot témoin le plus court. Posons $\{p,q\}$ comme étant cette paire et $w=a_1a_2\dots a_n$ le mot témoin le plus court qui les distingue. Encore une fois, un seul de $\hat{\delta}(p,w)$ et $\hat{\delta}(q,w)$ est acceptant.
	
	Ce mot $w$ ne peut pas être $\epsilon$. Auquel cas, la table aurait été remplie dès l'étape d'induction de l'algorithme.
	
	Ce mot $w$ doit forcément être de taille $\ge 1$ s'il n'est pas $\epsilon$. Considérons $r = \delta(p,a_1)$ et $s=\delta(q,a_1)$. Ces états sont différenciés par $a_2a_3\dots a_n$ puisque cette chaîne mène aux mêmes états que $\hat{\delta}(p,w)$ et $\hat{\delta}(q,w)$. Mais dans ce cas, cela signifie qu'il existe un mot plus petit que $w$ qui différencie deux états. Comme on a supposé que $w$ est le mot le plus petit qui différencie une mauvaise paire, $r$ et $s$ ne peuvent pas être une bad pair. Donc, l'algorithme a du découvrir qu'ils sont différentiables.
	
	Mais le pas d'induction stipule clairement que comme $\delta(p, a_1)$ et $\delta(q, a_1)$ mènent à deux états différentiables, $p$ et $q$ le sont aussi. On a une contradiction sur l'existence des mauvaises paires.
	
	Ainsi, s'il n'y en a pas, c'est que chaque paire différentiable est reconnue par l'algorithme.
\end{proof}


\subsubsection{Exemple}

\begin{figure}[H]
	\centering
	\begin{tikzpicture}[->,>=stealth',shorten >=1pt,auto,node distance=3cm, semithick, bend angle=10]
	
	\tikzstyle{every state}=[circle]
	
	\node[initial,state] (A)                    {$q_a$};
	\node[state]         (B) [below right of=A] {$q_b$};
	\node[state]         (C) [below left of=A] {$q_c$};
	\node[accepting,state]         (D) [below right of=B] {$q_d$};
	\node[state]         (E) [below left of=C]       {$q_e$};
	\node[state]         (F) [below right of=C]       {$q_f$};
	
	\path 	(A) 	edge              node {0} (C)
	edge              node {1} (B)
	(B) 	edge              node {0} (D)
	edge [bend left]  node {1} (F)
	(C) 	edge              node {0} (E)
	edge              node {1} (F)
	(D) 	edge [loop above] node {0,1} (D)
	(E) 	edge [loop above] node {0,1} (E)
	(F) 	edge              node {0} (D)
	edge [bend left]  node {1} (B);
	\end{tikzpicture}
	\caption{Automate $A_2$}\label{fig:a2}
\end{figure}

La première étape est de remplir la table avec l'algorithme précédant. Tout état est distinguable de $q_d$ : il est le seul état final. 5 cases peuvent déjà êtres cochées. Le reste de la table est remplie par induction.

\begin{figure}[H]
	\centering
	\begin{tabular}{ccccccc}
		\cline{2-2}
		\multicolumn{1}{c|}{B} & \multicolumn{1}{c|}{x} &&&&\\
		\cline{2-3}
		\multicolumn{1}{c|}{C} & \multicolumn{1}{c|}{x} &\multicolumn{1}{c|}{x}&&&\\
		\cline{2-4}
		\multicolumn{1}{c|}{D} & \multicolumn{1}{c|}{x} &\multicolumn{1}{c|}{x}&\multicolumn{1}{c|}{x}&&\\
		\cline{2-5}
		\multicolumn{1}{c|}{E} & \multicolumn{1}{c|}{x} &\multicolumn{1}{c|}{x}&\multicolumn{1}{c|}{x}&\multicolumn{1}{c|}{x}&\\
		\cline{2-6}
		\multicolumn{1}{c|}{F} & \multicolumn{1}{c|}{x} & \multicolumn{1}{c|}{}&\multicolumn{1}{c|}{x}&\multicolumn{1}{c|}{x}&\multicolumn{1}{c|}{x}\\
		\cline{2-6}
		\multicolumn{1}{c}{} & A&B&C&D&E\\
		
	\end{tabular}
	\caption{Table filling pour $A_2$, décelant des équivalences d'états}
	\label{fig:ta2}
\end{figure}


\subsubsection{Complexité}

Considérons $n$ le nombre d'états d'un automate, et $k$ la taille de l'alphabet $\Sigma$ supporté.

Si il y a $n$ états, il y a $\begin{pmatrix}n\\2\end{pmatrix}$ soit $\frac{n(n-1)}{2}$ paires d'états. A chaque itération (sur l'ensemble de la table), il faut considérer chaque paire, et vérifier si un de leur successeurs est différentiable. Cette étape prend au plus $O(k)$ pour tester chaque successeurs potentiel (en fonction du symbole lu).  Ainsi, une itération sur la table se fait en $O(kn^2)$. Si une itération ne découvre pas de nouveaux état différentiable s'arrête. Comme la table a une taille en $O(n^2)$ et qu'à chaque étape un élément au minimum doit y être coché, la complexité totale de l'algorithme est en $O(kn^4)$.

Cependant, il existe des pistes d'amélioration. La première est d'avoir, pour chaque paire $\{r,s\}$ une liste des paire $\{p,q\}$ qui, pour un même symbole, mènent à $\{r,s\}$. On dit de ces paires qu'elles sont dépendantes. Si la paire $\{r,s\}$ est marquée comme différenciable, leurs paires dépendantes seront de facto différenciables. 

Cette liste peut être construite en considérant chaque symbole $a \in \Sigma$ et ajoutant les paires $\{p,q\}$ à chacune de leur dépendance $\{\delta(p,a),\delta(q,a)\}$. Cette étape prend au plus $k.O(n^2)=O(kn^2)$. (Le nombre de symboles multiplié par le nombre de paires à considérer).

Ensuite, il suffit de partir des cas initiaux (se reposant sur le cas de base de l'algorithme), et de marquer tous leurs états dépendants comme différentiables, tout en ajoutant leur propre liste à chaque fois. La complexité de cette exploration est bornée par le nombre d'éléments dans une liste et le nombre de listes. Respectivement, $k$ et $O(n^2)$, ce qui donne $O(kn^2)$ pour cette exploration.

La complexité totale revient à $O(kn^2)$.




\subsection{Minimisation d'automate}\label{ss:miniauto}

La minimisation d'automate se fait en deux étapes :
\begin{enumerate}
	\item Se débarrasser de tous les états injoignables : ils ne participent pas à la construction du langage représenté
	\item Grâce aux équivalences d'états trouvées grâce au TFA déifni au point \ref{ss:tfa}, construire un nouvel automate. 
\end{enumerate}

Ces étapes vont être accompagnées d'un exemple, à savoir l'automate $A_1$ représenté à la figure \ref{fig:a1}.

\begin{figure}[H]
	\centering
	\begin{tikzpicture}[->,>=stealth',shorten >=1pt,auto,node distance=3cm, semithick, bend angle=10]
	
	\tikzstyle{every state}=[circle]
	
	\node[initial,state] (A)                    {$q_a$};
	\node[state]         (B) [below right of=A] {$q_b$};
	\node[state]         (C) [below left of=A] {$q_c$};
	\node[accepting, state]         (D) [below right of=B] {$q_d$};
	\node[state]         (E) [below left of=C]       {$q_e$};
	\node[state]         (G) [below right of=E]       {$q_g$};
	\node[state]         (F) [above right of=G]       {$q_f$};
	
	\path 	(A) 	edge              node {0} (C)
	edge              node {1} (B)
	(B) 	edge              node {0} (D)
	edge [bend left]  node {1} (F)
	(C) 	edge              node {0} (E)
	edge              node {1} (F)
	(D) 	edge [loop above] node {0,1} (D)
	(E) 	edge [loop above] node {0,1} (E)
	(F) 	edge              node {0} (D)
	edge [bend left]  node {1} (B)
	(G) 	edge              node {0} (E)
	edge              node {1} (F);
	\end{tikzpicture}
	\caption{Automate $A_1$}\label{fig:a1}
\end{figure}

L'état $q_g$ n'est pas atteignable : il peut être simplement supprimé. On obtient ainsi l'automate $A_2$ qui a servi d'exemple pour le TFA, représenté à la figure \ref{fig:a2}.

\subsubsection{Minimisation par table de différenciation}

Pour minimiser l'automate $A_2 = (Q, \Sigma, \delta, q_0, F)$, il faut :
\begin{enumerate}
	\item Générer la table de différenciation (qui, pour cet exemple, est à la figure \ref{fig:ta2})
	\item Séparer $Q$ en classes d'équivalences
	\item Construire l'automate canonique $A_3$:
		\begin{itemize}
			\item Soit $S$ une des classes d'équivalence
			\item Soit $\gamma$ la fonction de transition sur $S$. Pour un symbole $a \in \Sigma$, alors il doit exister une classe d'équivalence $T$ tel que pour chaque état $q$ dans $S$, $\delta(q,a) \in T$. Sinon, c'est que deux états $p$ et $q$ dans $S$ menant à différentes classes d'équivalences. Ces deux états sont différenciables, et ne pourraient pas appartenir tous deux à $S$ par construction. On peut écrire $\gamma(S,a)=T$.
		\end{itemize}
	\item L'état initial de $A_3$ est la classe d'équivalence contenant l'état initial de $A_2$ (dans notre exemple, l'état s'y trouve seul)
	\item Les états finaux ($F$) de $A_3$ sont les classes d'équivalences qui contenaient des états acceptants de $A_2$.
\end{enumerate}

La table de la figure \ref{fig:ta2}. Peut servir de base à la construction du nouvel automate suivant cet algorithme.
\begin{figure}[H]
	\centering
	\begin{tikzpicture}[->,>=stealth',shorten >=1pt,auto,node distance=3cm, semithick, bend angle=10]
	
	\tikzstyle{every state}=[circle]
	
	\node[initial,state] (A)                    {$q_a$};
	\node[state]         (B) [below right of=A] {$q_b$};
	\node[state]         (C) [below left of=A] {$q_c$};
	\node[accepting, state]         (D) [below right of=B] {$q_d$};
	\node[state]         (E) [below left of=C]       {$q_e$};
	
	\path 	
	(A) 	edge              node {0} (C)
	edge              node {1} (B)
	(B) 	edge              node {0} (D)
	edge [loop above] node {1} (B)
	(C) 	edge              node {0} (E)
	edge              node {1} (B)
	(D) 	edge [loop above] node {0,1} (D)
	(E) 	edge [loop above] node {0,1} (E);
	\end{tikzpicture}
	\caption{Automate $A_3$}
\end{figure}

Une expression régulière ($(1+01)1^*0(0+1)^*$) peut être déduite pour $L$ grâce à cet automate.


\subsection{Équivalence d'automates}

Considérons les automates $A_H$ et $A_I$ donnés dans les figures \ref{fig:ah} et \ref{fig:ai}

\begin{minipage}{0.4\linewidth}
	\begin{figure}[H]
		\centering
		\begin{tikzpicture}[->,>=stealth',shorten >=1pt,auto,node distance=2cm and 5cm, semithick, bend angle=10]
		
		\tikzstyle{every state}=[circle]
		
		\node[initial,state]	(A)					{$q_0$};
		\node[state]			(B)	[right= of A]	{$q_1$};
		\node[accepting,state]	(C) [below of=A]	{$q_2$};
		\node[accepting,state]	(D)	[below of=B]	{$q_3$};
		\node[accepting,state]	(E)	[below of=C]	{$q_4$};
		\node[state]			(F)	[below of=D]	{$q_5$};
		
		\path
		(A)	edge	[bend left]		node{0}		(B)
		(A)	edge					node{1}		(C)
		(B) edge	[bend left]		node{0}		(A)
		(B) edge					node{1}		(D)
		(C)	edge					node{0}		(E)
		(C)	edge					node[near start]{1}		(F)
		(D)	edge					node[near start, above]{0}		(E)
		(D)	edge					node{1}		(F)
		(E)	edge	[loop below]	node{0}	(E)
		(E) edge					node{1} (F)
		(F)	edge	[loop below]	node{0,1}	(F)
		
		; 
		\end{tikzpicture}
		\caption{Automate $A_H$, exemple d'un livre de référence\cite{Hopcroft79} (Fig3.2)}\label{fig:ah}
	\end{figure}
\end{minipage}\hspace{0.2\linewidth}
\begin{minipage}{0.4\linewidth}
	\begin{figure}[H]
		\centering
		\begin{tikzpicture}[->,>=stealth',shorten >=1pt,auto,node distance=1cm and 1cm, semithick, bend angle=10]
		
		\tikzstyle{every state}=[circle]
		
		\node[initial,state]	(A)					{$q_6$};
		\node[accepting,state]	(B)	[right= of A]	{$q_7$};
		\node[state]			(C) [right= of B]	{$q_8$};
		
		\path
		(A)	edge					node{1}		(B)
		(A)	edge	[loop above]	node{0}		(A)
		(B) edge					node{1}		(C)
		(B) edge	[loop above]	node{0}		(B)
		(C)	edge	[loop above]	node{0,1}	(C)
		
		; 
		\end{tikzpicture}
		\caption{Automate $A_I$, provenant également de \cite{Hopcroft79} }\label{fig:ai}
	\end{figure}
\end{minipage}

Ces deux automates sont-ils équivalents ? Grâce à l'algorithme de remplissage de table précédant, il est possible de tester l'équivalence d'états. Pour cela, il suffit de considérer les deux automates précédents comme un seul.

\begin{figure}[H]
	\centering
	\begin{tabular}{ccccccccc}
		\cline{2-2}
		\multicolumn{1}{c|}{$q_1$}&\multicolumn{1}{c|}{} &&&&&&&\\
		\cline{2-3}
		\multicolumn{1}{c|}{$q_2$}&\multicolumn{1}{c|}{x} &\multicolumn{1}{c|}{x}&&&&&&\\
		\cline{2-4}
		\multicolumn{1}{c|}{$q_3$}&\multicolumn{1}{c|}{x}&\multicolumn{1}{c|}{x}&\multicolumn{1}{c|}{}&&&&&\\
		\cline{2-5}
		\multicolumn{1}{c|}{$q_4$}&\multicolumn{1}{c|}{x}&\multicolumn{1}{c|}{x}&\multicolumn{1}{c|}{}&\multicolumn{1}{c|}{}&&&&\\
		\cline{2-6}
		\multicolumn{1}{c|}{$q_5$}&\multicolumn{1}{c|}{x}&\multicolumn{1}{c|}{x}&\multicolumn{1}{c|}{x}&\multicolumn{1}{c|}{x}&\multicolumn{1}{c|}{x}&&&\\
		\cline{2-7}
		\multicolumn{1}{c|}{$q_6$}&\multicolumn{1}{c|}{}&\multicolumn{1}{c|}{}&\multicolumn{1}{c|}{x}&\multicolumn{1}{c|}{x}&\multicolumn{1}{c|}{x}&\multicolumn{1}{c|}{x}&&\\
		\cline{2-8}
		\multicolumn{1}{c|}{$q_7$}&\multicolumn{1}{c|}{x}&\multicolumn{1}{c|}{x}&\multicolumn{1}{c|}{}&\multicolumn{1}{c|}{}&\multicolumn{1}{c|}{}&\multicolumn{1}{c|}{x}&\multicolumn{1}{c|}{x}&\\
		\cline{2-9}
		\multicolumn{1}{c|}{$q_8$}&\multicolumn{1}{c|}{x}&\multicolumn{1}{c|}{x}&\multicolumn{1}{c|}{x}&\multicolumn{1}{c|}{x}&\multicolumn{1}{c|}{x}&\multicolumn{1}{c|}{}&\multicolumn{1}{c|}{x}&\multicolumn{1}{c|}{x}\\
		\cline{2-9}
		\multicolumn{1}{c}{} & $q_0$& $q_1$ & $q_2$ & $q_3$ & $q_4$ & $q_5$ & $q_6$ & $q_7$\\
		
	\end{tabular}
	\caption{Table Filling, décelant des équivalences d'états entre les deux automates}\label{fig:tahi}
\end{figure}

De cette table, toujours grâce aux conclusions précédentes, il est possible d'extraire des classes d'équivalences : 
\begin{itemize}
	\item $C_0 = \{q_0, q_1, q_6\}$
	\item $C_1 = \{q_2, q_3, q_4, q_7\}$
	\item $C_2 = \{q_5, q_8\}$
\end{itemize}

En particulier, la classe $C_0$ souligne que les états initiaux sont équivalents. Cela signifie, par définition, que tout mot $w$ lu en partant de cet état sera soit accepté dans les deux automates, soit refusé dans les deux. $A_H$ et $A_I$ définissent donc le même langage.



\subsection{Construction d'automate depuis un langage}

Soit le langage $A_N = \{w | w \text{ fini par b et ne contient pas bb}\}$ défini sur $\Sigma_N = {a,b}$.

On peut diviser les mots en 3 ensembles : 

\begin{itemize}
	\item $W_0$ le sous-ensemble des mots ne finissant pas le symbole $b$
	\item $W_1$ celui des mots finissant par le symbole $b$ mais ne contenant pas $bb$
	\item $W_2$ celui des mots contenant au moins $bb$
\end{itemize}

Il y a d'autres façons de construire des sous-ensembles, mais celle-ci à l'avantage de rendre la question de l'appartenance à $L_N$ triviale : un mot appartient au second ensemble si et seulement si il fait partie du langage, par définition.

De plus, tous les éléments d'un sous-ensemble respectent la relation $R_L$ entre eux. ($R_L : xR_Ly \Leftrightarrow \forall z \in \Sigma^*, xz \in L \Leftrightarrow yz \in L$). Cela en fait des classes d'équivalence sur cette relation.

Cela peut être démontré pour chaque sous-ensemble :
\begin{itemize}
	\item Soient $x,y \in W_0$. Soit $z \in \Sigma^*$. Dès lors, si $xz \in L_N$, c'est que $z$ fini par $b$ mais ne contient pas $bb$, et donc $yz \in L_N$. Si $yz \in L_N$, le même argument peut être appliqué.
	\item Soient $x,y \in W_1$. Soit $z \in \Sigma^*$. Dès lors, si $xz \in L_N$, c'est que $z$ ne commençait pas le symbole $b$ et ne contenait pas $bb$, $yz$ ne contiendra donc pas $bb$, puisque cette chaîne n'est ni dans $z$ ni dans $y$, ni a cheval sur les deux, $z$ ne commençant pas par $b$. Ainsi, $yz \in L_N$. Si $yz \in L_N$, le même argument peut être appliqué.
	\item Soient $x,y \in W_2$. Soit $z \in \Sigma^*$. Comme $x$ contient déjà $bb$, $x \notin L_N$ et, a fortiori, $xz \notin L_N$. Comme la prémisse est fausse, l'implication $xz \in L \Rightarrow yz \in L$ est vraie. La même logique peut être appliquée à partir de $y$ pour justifier l'implication inverse.
\end{itemize}

De plus, ces sous-ensembles sont disjoints. Cela peut se prouver en invalidant la relation pour certains éléments entre eux, mais dans ce cas-ci, la propriété est assurée par définition.

Ceci revient à démontrer que $W_0,W_1,W_2$ sont des classes d'équivalence. De plus, $R_L$ respecte la congruence à droite, comme démontré dans la preuve du théorème de Myhill-Nérode. Ce même théorème donne une méthode pour construire un automate : prendre un représentant pour chaque classe et en faire un état.

\begin{itemize}
	\item $\Sigma=\{a,b\}$ est connu.
	\item $Q=\{[[\epsilon]]\, [[b]], [[bb]]\} = \{q_\epsilon, q_b, q_{bb}\}$
	\item $q_0 = q_\epsilon$ 
	\item $F = \{q_b\}$ l'union des classes acceptant
	\item $\delta$ défini en utilisant des exemples tirés des classes d'équivalence.
\end{itemize}

Ce qui donne l'automate de la figure \ref{fig:an}

\begin{figure}[H]
	\centering
	\begin{tikzpicture}[->,>=stealth',shorten >=1pt,auto,node distance=2cm and 2cm, semithick, bend angle=10]
	
	\tikzstyle{every state}=[circle]
	
	\node[initial,state]	(A)					{$q_\epsilon$};
	\node[accepting,state]	(B)	[right= of A]	{$q_b$};
	\node[state]			(C) [right= of B]	{$q_{bb}$};
	
	\path
	(A)	edge	[bend left]		node{b}		(B)
	(A)	edge	[loop above]	node{a}		(A)
	(B) edge	[bend left]		node{a}		(A)
	(B) edge					node{b}		(C)
	(C)	edge	[loop above]	node{a,b}	(C)
	
	; 
	\end{tikzpicture}
	\caption{Automate $A_N$, exemple d'une thèse\cite{Neider14}}\label{fig:an}
\end{figure}

Cet automate est bien une représentation du langage $L_N$. Seul un mot finissant par $b$ mais ne contenant pas $bb$ se termine à l'état $q_b$.



	
	
	\section{Preuves}\label{sec:preuves}Cette section apporte le détail sur la relation de Myhill-Nérode, en prouve les propriétés avant d'en faire l'usage dans le théorème du même nom.




\subsection{Relation de Myhill-Nérode}

Soit un langage $L$ sur un alphabet $\Sigma$.

Soit la relation $R_L$ portant sur deux mots (ne faisant pas nécessairement partie de $L$). Deux mots $x$ et $y$ respectent la relation de Myhill-Nérode $R_L$ si 

$$\forall z \in \Sigma^*, xz \in L \Leftrightarrow yz \in L$$

Intuitivement, deux mots sont en relation si pour tout mot qu'on leur concatène, les deux mots résultants sont tous deux dans le langage ou non.

\begin{lemma}
	Cette relation est une relation d'équivalence. De plus, elle respecte la congruence à droite. C'est à dire que si $xR_Ly$, alors pour tout symbole $a \in \Sigma$, $xaR_Lya$
\end{lemma}

\begin{proof}[Equivalence et Congruence à droite]
	Dire d'une relation qu'elle décrit une équivalence, revient à dire qu'elle est réflexive, transitive et symétrique
\begin{itemize}
		\item $R_L$ est réflexive. Soit $x \in \Sigma^*$. Soit $z \in \Sigma^*$. Montrer que $xR_Lx$ est vrai revient à montrer que $ xz \in L \Leftrightarrow xz \in L$ est vrai. $R_L$ est donc réflexive.
		\item $R_L$ est symétrique. Soient $x, y \in \Sigma^*$ tels que $xR_Ly$. Soit $w \in \Sigma^*$. Montrer que $yR_Lx$ revient à montrer que $ yw \in L \Leftrightarrow xw \in L$. Or, par hypothèse, $ xz \in L \Leftrightarrow yz \in L$, qui peut s'écrire aussi $ yz \in L \Leftrightarrow xz \in L$ pour tout $z \in \Sigma^*$, et en particulier $z=w$.
		\item $R_L$ est transitive. Soient $x,y,u \in \Sigma^*$ tels que $xR_Ly$ et $yR_Lz$. Soit $w \in \Sigma^*$. Comme $ xz \in L \Leftrightarrow yz \in L$ et $ yz \in L \Leftrightarrow uz \in L$ pour tout $z \in \Sigma^*$ (par hypothèse), c'est vrai en particulier pour $z=w$. Dès lors,  $ xw \in L \Leftrightarrow yw \in L$et $ yw \in L \Leftrightarrow uw \in L$. Par transitivité de l'implication, on obtient $ xw \in L \Leftrightarrow uw \in L$, à savoir $xR_Lu$.
		\item $R_L$ est congruente à droite. Soient $x,y \in \Sigma^*$ tels que $xR_Ly$. Soit $a \in \Sigma$. Par hypothèse, $ xz \in L \Leftrightarrow yz \in L$ pour tout $z \in \Sigma^*$. Cela doit donc être vrai en particulier pour le mot $z=aw$ avec $w$ quelconque. En remplaçant dans l'hypothèse, on obtient  $ xaw \in L \Leftrightarrow yaw \in L$. Ce qui montre que $xaR_Lya$.
	\end{itemize}
\end{proof}


\subsection{Théorème de Myhill-Nerode}
	
	\begin{theorem}
		Les 3 énoncés suivants sont équivalents :
		\begin{enumerate}
			\item Un langage $L\subseteq\Sigma^*$ est accepté par un DFA
			\item $L$ est l'union de certaines classes d'équivalence d'index fini respectant une relation d'équivalence et de congruence à droite
			\item Soit la relation d'équivalence $R_L : xR_Ly \Leftrightarrow \forall z \in \Sigma^*, xz \in L \Leftrightarrow yz \in L$ (la relation de Myhill-Nérode définie précédemment). $R_L$ est d'index fini.
		\end{enumerate}
	\end{theorem}
	
	\begin{proof}La preuve d'équivalence se fait en prouvant chaque implication de façon cyclique :\\
		
		$(1)\rightarrow(2)$ Supposons que (1) soit vrai, c'est à dire que le langage $L$ est accepté par un automate déterministe fini $A$. Considérons la relation d'équivalence $R_M$ étant vraie pour les mots $x,y$ si $\hat{\delta}(q_0,x)\in F \iff \hat{\delta}(q_0,y)\in F$. Elle a été définie en \ref{ss:rm}. Il y est prouvé qu'elle est congruente à droite. Comme il y a au plus une classe d'équivalence pour $R_M$ par état de $A$. Comme ce nombre d'états est fini, $R_M$ est d'index fini. De plus, $L$ est l'union de classes contenant un mot $w$ tel que $\hat{\delta}(q_0,w) \in F$, (or, ce chemin retourne un état. Il s'agit donc d'une union des classes correspondant aux états acceptants).
		
		$(2)\rightarrow(3)$ Montrons que pour toute relation $E$ satisfaisant (2), chaque classe est intégralement contenue dans une seule classe de $R_L$. Ces classes étant d'index fini, c'est un argument suffisant pour déduire que $R_L$ est d'index fini. Considérons $x,y$ tels que $xEy$. Comme $E$ est congruente à droite, pour tout mot $z \in \Sigma^*$, on sait que $xzEyz$. Comme $L$ est un union de ces classes d'équivalence, $xzEyz$ implique que $xz \in L \Leftrightarrow yz \in L$, ce qui revient à $xR_Ly$. Cela signifie que tout mot dans la classe d'équivalence de $x$ définie par $E$ se retrouve dans la même classe d'équivalence que $x$ par $R_L$. Ce qui permet de conclure que chaque classe d'équivalence de $E$ est contenue dans une classe d'équivalence de $R_L$. 
	
		
		$(3)\rightarrow(1)$ Considérons la relation $R_L$ définie précédemment, et déduisons-en $Q'$ les classes d'équivalence sur $L$ et $[[x]]$ l'élément(la classe) de $Q'$ qui contient $x$. Puisque $R_L$ a été démontré comme congruent à droite, on peut définir des transitions : $\delta'([[x]],a) = [[xa]]$. En choisissant un élément $y$ dans $[[x]]$ (ce qui signifie que $xR_Ly$), on obtient $\delta'([[x]],a)=[[ya]]$. Sauf que par définition, $xR_Ly$ signifie qu'en y ajoutant n'importe quel mot $z$, $xz$ et $yz$ appartiennent tous deux où non à $L$. C'est vrai en particulier pour $z=az'$. Ainsi, $xaz'$ et $yaz'$ appartiennent tous deux à $L$ ou non. Ce qui signifie que $xaR_Lya$ et donc $[[xa]]=[[ya]]$. Posons $q_0'=[[\epsilon]]$ et $F' = \{[[x]]|x \in L\}$. Tous ces éléments forment l'automate $M'=(Q', \Sigma, \delta', q_0', F')$. Il est déterministe par la définition de $\delta'$, fini car $Q'$ est fini par construction (le nombre de classes d'équivalence est fini). De plus, il accepte $L$ puisque $\delta'(q_0',x)=[[x]]$, ce qui signifie que $x \in L(M')$ si et seulement si $[[x]] \in F'$, qui a été défini comme tel.
		
		
	\end{proof}


	\begin{corollary}
		Grâce à la preuve du théorème de Myhill-Nérode, en particulier la justification partant de la relation d'équivalence $R_L$ pour montrer que la langage $L\subseteq\Sigma^*$ est accepté par un DFA, on a une méthode pour construire un automate à partir d'un langage.
	\end{corollary}

	
	\newpage
	\bibliographystyle{siam}
	\bibliography{refs.bib}
	
\end{document}